% \documentclass[a4paper,12pt]{article}{{{
\documentclass[a4paper,11pt]{article}
% -------------------------------------------------
% Set up the page margins.
% -------------------------------------------------
% \usepackage[text={6.5in,9.5in},centering]{geometry}
% \setlength{\topmargin}{-0.25in}
\usepackage[margin=1in]{geometry}

% -------------------------------------------------\left( \left[ \frac{a-d}{2}\right]q + p \right)
% Load Packages here
% -------------------------------------------------
%\usepackage{hyperref}
\usepackage{graphicx,url,subfig}
\RequirePackage[OT1]{fontenc}
\RequirePackage{amsthm,amsmath,amssymb,amscd}
\RequirePackage[numbers]{natbib}
\RequirePackage[colorlinks,citecolor=blue,urlcolor=blue]{hyperref}
\usepackage{graphics}
\usepackage{enumerate}
\usepackage{datetime}
\usepackage{etex}
% \usepackage[notcite,notref,color]{showkeys}
\usepackage[normalem]{ulem} % either use this (simple) or
\usepackage{soul} % use this (many fancier options)
\makeatother
\numberwithin{equation}{section}
\allowdisplaybreaks
\usepackage{mathrsfs}


% -------------------------------------------------
% check references
% -------------------------------------------------
% \usepackage{refcheck}
% \usepackage[notref,notcite]{showkeys}
% \usepackage[notcite]{showkeys}
% \usepackage{showkeys}

% -------------------------------------------------
% Environments here
% -------------------------------------------------
\theoremstyle{plain}
\newtheorem{theorem}{Theorem}[section]
\newtheorem{corollary}[theorem]{Corollary}
\newtheorem{lemma}[theorem]{Lemma}
\newtheorem{proposition}[theorem]{Proposition}
\newtheorem{exercise}{Exercise}
\theoremstyle{definition}
\newtheorem{definition}[theorem]{Definition}
\newtheorem{hypothesis}[theorem]{Hypothesis}
\newtheorem{assumption}[theorem]{Assumption}

\newtheorem{qu}{Question}

\newtheorem{remark}[theorem]{Remark}
\newtheorem{conjecture}[theorem]{Conjecture}
\newtheorem{example}[theorem]{Example}

% -------------------------------------------------
%  Define short hand symbols.
% -------------------------------------------------
\newcommand{\E}{\mathbb{E}}
\newcommand{\W}{\dot{W}}
\newcommand{\B}{\textbf{B}}
\newcommand{\D}{\mathbb{D}}
\newcommand{\ud}{\ensuremath{\mathrm{d} }}
\newcommand{\Ceil}[1]{\left\lceil #1 \right\rceil}
\newcommand{\Floor}[1]{\left\lfloor #1 \right\rfloor}
\newcommand{\sgn}{\text{sgn}}
\newcommand{\Lad}{\text{L}_{\text{ad}}^2}
\newcommand{\SI}[1]{\mathcal{I}\left[#1 \right]}
\newcommand{\SIB}[2]{\mathcal{I}_{#2}\left[#1 \right]}
\newcommand{\Indt}[1]{1_{\left\{#1 \right\} }}
\newcommand{\LadInPrd}[1]{\left\langle #1 \right\rangle_{\text{L}_\text{ad}^2}}
\newcommand{\LadNorm}[1]{\left|\left|  #1 \right|\right|_{\text{L}_\text{ad}^2}}
\newcommand{\Norm}[1]{\left\|  #1   \right\|}
\newcommand{\Ito}{It\^{o} }
\newcommand{\Itos}{It\^{o}'s }
\newcommand{\spt}[1]{\text{supp}\left(#1\right)}
\newcommand{\InPrd}[1]{\left\langle #1 \right\rangle}
\newcommand{\mr}{\textbf{r}}
\newcommand{\Ei}{\text{Ei}}
\newcommand{\arctanh}{\operatorname{arctanh}}
\newcommand{\ind}[1]{\mathbb{I}_{\left\{ {#1} \right\} }}

\DeclareMathOperator{\esssup}{\ensuremath{ess\,sup}}

\newcommand{\steps}[1]{\vskip 0.3cm \textbf{#1}}
\newcommand{\calB}{\mathcal{B}}
\newcommand{\calD}{\mathcal{D}}
\newcommand{\calE}{\mathcal{E}}
\newcommand{\calF}{\mathcal{F}}
\newcommand{\calG}{\mathcal{G}}
\newcommand{\calK}{\mathcal{K}}
\newcommand{\calH}{\mathcal{H}}
\newcommand{\calI}{\mathcal{I}}
\newcommand{\calL}{\mathcal{L}}
\newcommand{\calM}{\mathcal{M}}
\newcommand{\calN}{\mathcal{N}}
\newcommand{\calO}{\mathcal{O}}
\newcommand{\calT}{\mathcal{T}}
\newcommand{\calP}{\mathcal{P}}
\newcommand{\calR}{\mathcal{R}}
\newcommand{\calS}{\mathcal{S}}
\newcommand{\bbC}{\mathbb{C}}
\newcommand{\bbN}{\mathbb{N}}
\newcommand{\bbP}{\mathbb{P}}
\newcommand{\bbZ}{\mathbb{Z}}
\newcommand{\myVec}[1]{\overrightarrow{#1}}
\newcommand{\sincos}{\begin{array}{c} \cos \\ \sin \end{array}\!\!}
\newcommand{\CvBc}[1]{\left\{\:#1\:\right\}}
\newcommand*{\one}{ {{\rm 1\mkern-1.5mu}\!{\rm I} }}
\newcommand{\RR}{\mathbb{R}}
\newcommand{\R}{\mathbb{R}}
\newcommand{\myEnd}{\hfill$\square$}
\newcommand{\ds}{\displaystyle}
\newcommand{\Shi}{\text{Shi}}
\newcommand{\Chi}{\text{Chi}}
\newcommand{\Daw}{\ensuremath{\mathrm{daw} }}
\newcommand{\Erf}{\ensuremath{\mathrm{erf} }}
\newcommand{\Erfi}{\ensuremath{\mathrm{erfi} }}
\newcommand{\Erfc}{\ensuremath{\mathrm{erfc} }}
\newcommand{\He}{\ensuremath{\mathrm{He} }}

\def\crtL{\rho}

\newcommand{\conRef}[1]{(#1)}

\DeclareMathOperator{\Lip}{\mathit{L}}
\DeclareMathOperator{\LIP}{Lip}
\DeclareMathOperator{\lip}{\mathit{l}}

\DeclareMathOperator{\Vip}{\overline{\varsigma}}
\DeclareMathOperator{\vip}{\underline{\varsigma}}
\DeclareMathOperator{\vv}{\varsigma}

% % Set up mdframe for questions.
\newcounter{NQ}
\newcounter{LQ}
\newcounter{ZQ}
\usepackage{xcolor}
\usepackage[framemethod=tikz]{mdframed}
\mdfsetup{%
	middlelinecolor=lightgray,
	middlelinewidth=2pt,
	backgroundcolor=red!20,
	roundcorner=10pt}
\newenvironment{NickQuestion}[1][] %
{\refstepcounter{NQ}
	\begin{mdframed}[backgroundcolor=lightgray]\textbf{Nick's Question \theNQ \: #1~~}~~~\noindent}%
	{\end{mdframed}}
\newenvironment{LeQuestion}[1][] %
{\refstepcounter{LQ}
	\begin{mdframed}[backgroundcolor=red!10]\textbf{Le's Question \theLQ \: #1~~}~~~\noindent}%
	{\end{mdframed}}
\newenvironment{ZhijianQuestion}[1][] %
{\refstepcounter{ZQ}
	\begin{mdframed}[backgroundcolor=blue!10]\textbf{Zhijian's Question \theZQ \: #1~~}~~~\noindent}%
	{\end{mdframed}}
\numberwithin{NQ}{section}
\numberwithin{LQ}{section}
\numberwithin{ZQ}{section}

%%%%%%%%%%%%%%%%%%%%%
\usepackage[backend=biber,
style=alphabetic,
hyperref=true,
backref=true,
natbib=true,
sorting=nyt,
abbreviate=true]{biblatex}
\AtEveryBibitem{% Clean up the bibtex rather than editing it
	\clearfield{doi}
	\clearfield{url}
	\clearfield{issn}
	\clearlist{location}
	\clearfield{month}
	\clearfield{series}
	
	\ifentrytype{book}{}{% Remove publisher and editor except for books
		\clearlist{publisher}
		\clearname{editor}
	}
}
\addbibresource{All.bib}
%%%%%%%%%%%%%%%%%%%%%%%

\usepackage[notref,notcite]{showkeys}

\title{Nick Eisenberg's Meeting Notes}
\author{Le Chen, Nick Eisenberg, Zhijian Wu}
\date{October 2019}

\usepackage{natbib}
\usepackage{graphicx}
% }}}


\title{Global solutions for the interpolated stochastic heat and wave equation with a super-linear diffusion term}
\author{Le Chen\\Auburn University\\\texttt{le.chen@auburn.edu}\\
	\and
	Nicholas Eisenberg\\Auburn University\\\texttt{nze0019@auburn.edu}
}
\date{\today}

\begin{document}
\maketitle

\begin{abstract}
In this paper we prove the existence of a global solution to the interpolated stochastic heat and wave equation driven by a space-time white noise and with a super-linear non-linear diffusion term.
\end{abstract}
\bigskip

\noindent{\it \noindent MSC 2010 subject classification}:
\smallskip

\noindent{\it Keywords}: 
\smallskip

\tableofcontents
	\section{Introduction}
	In this chapter, we study the interpolated stochastic heat and wave equation (ISHWE) on the whole space $\R^d$,
	\begin{equation} \label{E:SupLinSPDE}
	\begin{cases}
	\left(\partial^\beta_t + \frac{\nu}{2}(-\Delta)^{\alpha/2}\right)\: u(t,x) = I^\gamma_t \left[\sigma( u(t,x)) \: \dot W(t,x) \right] & \text{$x\in \R^d$, $t>0$} \\
	u(0,\cdot) = u_0(x)                                                                                                         & \beta \in (0,1]               \\
	u(0,\cdot) = u_0(x), \quad \partial_t u(0,\cdot) = v_0(x)                                                                            & \beta \in (1,2),
	\end{cases}
	\end{equation}
	where the fractional differential operators $\partial^\beta_t$, $(-\Delta)^{\alpha/2}$ and $I^\gamma_t$ respectively denote the {\em Caputo derivative}, the {\em fractional Laplacian} and the {\em Riemann-Liouville fractional integral operator}. The noise, $\dot{W}$ is taken to be space-time white and the initial data, $u_0$ and $v_0$, are assumed to be Borel measurable functions which are Bounded on $\R^d$. The nonlinearity, $\sigma$, is {\em locally Lipschitz} in the sense that the following asymptotic relation holds as $|x|,|z| \to \infty$:
	\begin{equation}
	\label{C:SigSupLin}
	|\sigma(x)-\sigma(z)| \le \sigma_2|x-z|[\ln_+(|x-z|)]^{\delta},
	\end{equation}
	where $\sigma_2 , \delta >0$ and $\ln_+(z) = \ln(z \vee e)$ for $z >0$. Further, this implies that as $|x| \to \infty$,
	\begin{equation}
	|\sigma(x)| \le \sigma_1 + \sigma_2|x|[\ln(|x|)]^{\delta},
	\end{equation}
	where $\sigma_1 = \sigma(0)$.
	
For any $T>0$ with $t \in [0,T]$ and $x \in \R^d$, the solution to \eqref{E:SupLinSPDE} is understood in the {\em mild} sense as the solution to the following stochastic integral equation:
\[
u(t,x) = J_0(t,x) + I(t,x),
\]
	where
	\begin{align*}
	J_0(t,x) &=\begin{cases} 
	[Z(t,\cdot)*u_0](x) & \beta \in (0,1]
	\\	[Z(t,\cdot)*v_0](x) + [Z^*(t,\cdot)*u_0](x) & \beta \in (1,2)
	\end{cases}
	\\	I(t,x) &= \int_0^t \int_{\R^d} Y(t-s,x-y)\sigma(u(s,y)) W(\ud s, \ud y).
	\end{align*}

The stochastic integral above is a {\em Walsh integral} (e.g. see \cite{walsh:86:introduction}) and the fundamental solution is a triple, $\{Z, Z^*, Y\}$, and each member of the triple is defined through the Fox-H function, however, one can more compactly define them through their Fourier transforms \cite[Theorem 4.1]{chen.hu.ea:19:nonlinear}: 
	\begin{align*}
		\mathcal{F}Z(t,\cdot)(\xi) &= t^{\lceil \beta \rceil -1} E_{\beta, \lceil \beta \rceil}\left(-2^{-1} \nu t^\beta |\xi|^{\alpha}\right), \\
		\mathcal{F}Z^*(t,\cdot)(\xi) &= E_{\beta, 1}\left(-2^{-1} \nu t^\beta |\xi|^{\alpha}\right), \\
		\mathcal{F}Y(t,\cdot)(\xi) &= t^{\beta + \gamma-1}E_{\beta. \beta +\gamma}\left(-2^{-1}\nu t^\beta |\xi|^\alpha\right),
	\end{align*}
where $E_{a, b}(z)$ denotes the {\em two-parameter Mittag-Leffler function} \cite{podlubny:99:fractional}:
	\[
		E_{a,b}(z) := \sum_{k=0}^{\infty} \frac{z^k}{\Gamma(ak+b)}, \quad a>0, \: b>0,
	\]
and $\Gamma$ is the standard {\em Gamma-function}. 

We mention that under the special case where $\alpha \in (1,2)$, $\beta \in (0,1]$ and $\gamma = 1 - \beta$, that $Z = Z^* = Y$, which is easily seen from their Fourier transforms given above. Because of these equalities, we simply let $Y$ denote the fundamental solution under this situation. Moreover, under this case there is a probabilistic representation for $Y$ which is used by the authors in \cite{foondun.liu.ea:19:some, foondun.nane:17:asymptotic, mijena.nane:15:space-time, foondun.liu.ea:17:moment, mijena.nane:16:intermittence}.  Indeed, let $X_t$ denote a symmetric $\alpha$-stable process with density $p(t,x)$. Let $D=\{D_r, \: r \ge 0\}$ denote a $\beta$-stable subordinator and $E_t$ its first passage time. Then it is known that the density of the time changed processes $X_{E_t}$ is given by $Y(t,x)$ and we have that 
	\[
		Y(t,x) = \int_0^\infty p(s,x)f_{E_t}(s) \: \ud s,
	\]
where
	\[
		f_{E_t}(s) = t \beta^{-1} x^{-1-1/\beta}g_{\beta}\left(tx^{-1/\beta} \right),
	\]
where $g_\beta(\cdot)$ is the density function of $D_1$ and is smooth on the real line with $g_{\beta}(x) = 0$ for $x \le 0$.

The goal of this paper is to prove the existence and uniqueness of a global solution to \eqref{E:SupLinSPDE}. The work is highly motivated by the recent work by Millet et al \cite{millet.sanz-sole:21:global}. It is a well studied phenomena that either a super-linear drift or diffusion term may cause blow-up of the solution. As for the stochastic heat equation, we direct the reader to \cite{foondun.nualart:21:osgood, mueller.sowers:93:blowup, dalang.khoshnevisan.ea:19:global, bonder.groisman:09:global}. On the other hand, the only other work that we are aware of that is dedicated to proving some non-existence results of \eqref{E:SupLinSPDE} is \cite{asogwa.mijena.ea:20:blow-up}. So, to the best of our knowledge, this is the first work on proving the global existence of a solution to \eqref{E:SupLinSPDE} with a super-linear diffusion term (see Theorem \ref{T:MainSupLin} below). However, as we will see, the calculations given below are essentially identical to those performed in \cite{millet.sanz-sole:21:global} which is the cause for our motivation. \bigskip 

Here is the main result of this chapter. The proof will be postponed until Section \ref{S:TMainSupLin} below.

	\begin{theorem} \label{T:MainSupLin}
	Suppose that $\sigma$ satisfies \eqref{C:SigSupLin} and the initial data are Borel measurable functions. Then for any $R>0$ and $(t,x) \in [0,T] \times \bar{B}_R(0)$, where $\bar{B}_R(0)$ denotes the closed ball centered around the origin, there exists a random field solution to \eqref{E:SupLinSPDE} denoted as $(u(t,x): (t,x) \in [0,T] \times \bar{B}_R(0))$. This solution is unique and satisfies 
	\begin{equation}
	\sup_{t\in[0,T],\; |x| \le R} |u(t,x)| < \infty, \quad \text{almost surely.}
	\end{equation}
\end{theorem}

The proof of Theorem \ref{T:MainSupLin} uses a standard stopping time argument along with a series of truncations of the super-linear term, $\sigma$, in order to construct a solution of \eqref{E:SupLinSPDE}. Indeed, we define 	
\[ \sigma_N(x) = \begin{cases}
\sigma(x) & |x| \le N \\
\text{continuous and constant} & |x| > N,
	\end{cases}
\]
and then associate for each $N>0$, the corresponding global solution $u_N(t,x)$, whose existence follows due to Theorem \ref{T:SolveabilityGenLip} below. We then define the stopping time $\tau_N$ as follows:
\[
\tau_N := \inf \left\{ t>0 : \sup_{|x| \le R} |u_N(t,x)| \ge N \right\} \wedge T.
\]
We prove that $\{t < \tau_N\} \uparrow \Omega$ and then show that the solution to \eqref{E:SupLinSPDE} exists pathwise for $\omega \in \Omega$. We remind the reader that in the case of a globally Lipschitz diffusion term, one usually proves that the solution exists as a limit in $L^2(\Omega)$. Because of the stopping-time argument, existence in $L^2(\Omega)$ is no longer guaranteed. Moreover, because $L^2(\Omega)$ existence is no longer guaranteed, then the solution may no longer satisfy the standard Ito isometry. 
	
	\section{Preliminary results}
	
	In this section, we assume that the coefficient $\sigma$ is globally Lipschitz and therefore satisfies the following inequality:
	\begin{equation} \label{E:SupLinGlobalLip}
	|\sigma(x)| \le c(\sigma) + L(\sigma)|x|,
	\end{equation}
	where $c(\sigma) = \sigma(0)$ and $L(\sigma)$ is the Lipschitz constant of $\sigma$ an we assume that $L(\sigma)>0$. The equation \eqref{E:SupLinSPDE} under the assumption that $\sigma$ satisfies \eqref{E:SupLinGlobalLip} has been studied in \cite{chen.hu.ea:19:nonlinear} where they establish in Theorem 4.1 {\em ibid} the existence and uniqueness of global random field solutions. In order to do so, one needs to show {\em Dalang's condition}:
	\begin{equation}
	\label{C:Dalang}
	\int_0^t \ud s \int_{\R^d} \ud y |Y(s,y)|^2 < \infty, \quad \text{for all } t>0,
	\end{equation} 
	which is equivalent to (e.g. see Lemma 5.3 {\em ibid})
	\begin{equation}
	\label{E:DalangEquiv}
	d < 2\alpha + \frac{\alpha}{\beta}\min\{2\gamma -1, 0\} =: \Theta,
	\end{equation}
	which is further equivalent to 
	\begin{equation} \label{E:DalangEquiv2}
	\rho(d) > 0 \quad \text{and} \quad d < 2\alpha,
	\end{equation}
	where 
	\begin{equation} \label{D:rho}
	\rho(x) := 2\beta + 2\gamma -1 -\beta x / \alpha.
	\end{equation} 

We now state the existence and uniqueness result proven in \cite{chen.hu.ea:19:nonlinear}.
	
	\begin{theorem}{\cite[Theorem 4.1]{chen.hu.ea:19:nonlinear}}
		\label{T:SolveabilityGenLip}
		Under \eqref{C:Dalang}, the SPDE \eqref{E:SupLinSPDE} has a unique (in the sense of versions) random field solution $\{u(t,x) : (t,x) \in (0,\infty) \times \R^d \}$ if the initial data are such that 
		\begin{equation} \label{E:HomogBdd}
		\widehat{C}_T := \sup_{t\in[0,T],\; x\in \R^d} |J_0(t,x)| < \infty.
		\end{equation}	
%		\textcolor{blue}{Moreover, the following statements are true:
%			\begin{enumerate}
%				\item $(t,x) \mapsto u(t,x)$ is $L^P(\Omega)$-continuous for all $p \ge 2$,
%				\item For all even integers $p \ge 2$, all $t \ge 0$ and $x,y\in\R^d$,
%				\begin{equation}
%				\Norm{u(t,x)}_p^2 \le 2J_0^2(t,x) + \left[ \bar{\varsigma}^2 + 2 \hat{C}_t^2 \right]\exp(C \lambda^{2/(1-\sigma)}t).
%				\end{equation}
%		\end{enumerate}}
		
	\end{theorem}

We now state a lemma proven in \cite{chen.hu.ea:19:nonlinear} that calculates the $L^1(\R^d)$ and $L^2(\R^d)$ norms of the fundamental solutions. The lemma will be needed below in some of our calculations.

	\begin{lemma}{\cite[Theorem 4.1 and Lemma 5.5]{chen.hu.ea:19:nonlinear}}
		\label{L:1and2normYZ}
		Assume that $d < 2\alpha$, $\beta \in (0,2)$, and $\gamma \ge 0$. Then 
		\[
		\int_{\R^d} Y(t,x) \ud x = t^{\beta+\gamma -1 } \quad \text{and} \quad \int_{\R^d} Z(t,x) \ud x= t^{\lceil \beta \rceil -1}
		\]
		and
		\[
		\int_{\R^d} Y^2(t,x) \ud x = C_{\#} t^{2(\beta+\gamma-1)-d\beta/\alpha} = C_{\#} t^{\rho -1} ,
		\]
		for all $t > 0$, where $\rho = \rho(d)$ and 
		\[
		C_{\#} := \frac{2}{\Gamma(d/2)(4 \pi )^{d/2}} \int_0^\infty u^{d-1}E^2_{\beta,\beta + \gamma}(-u^\alpha) \ud u.
		\]
		Now further assume that $\beta \in (1,2)$, then 
		\[
		\int_{\R^d} Z^*(t,x) \ud x= 1. 
		\]
	\end{lemma}
	\begin{remark}
		When $\alpha = \beta = 2$,  $\gamma = 0$ and $d=1$ then Lemma \ref{L:1and2normYZ} gives us that $\int_{\R} Y^2(t,x) \ud x = t/2$ which coincides with the $L^2(\R)$ norm of the wave kernel. 
	\end{remark}

\subsection{Some moment bounds}

In this section, we will prove some moment bounds and a continuity result for the solution with a globally Lipschitz diffusion term. In the proof of the next proposition, we will use the following two facts:
	\begin{equation}
	\label{E:SUPidentities}
	\sup_{t\ge0} t^ke^{-at} = k^k(ea)^{-k} \text{ for }k \in \mathbb{N} \text{ and } \sup_{t \ge 0} \int_0^t s e^{-as} \ud s = a^{-2} \text{ for } a>0.
	\end{equation}
	
	\begin{proposition}{\normalfont\cite[Proposition 3.2]{millet.sanz-sole:21:global}}
		\label{P:N(u)Bdd}
		Let $u_0$ and $v_0$ be Borel functions satisfying $\Norm{u_0}_\infty + \Norm{v_0}_\infty < \infty$. Suppose that $d < 2 \alpha$ and $\rho = \rho(d) > 0$, where $\rho(x)$ is defined in \eqref{D:rho}. Then there exists a universal constant $K:=K(\alpha, \beta, \gamma,d) >0$ such that for any 
		\[
		a > \left( 8 L^2(\sigma)K^2\right)^{1/\rho}.
		\]
		and 
		\[
		p \in \left[ 2, \frac{a^{\rho}}{4L^2(\sigma)K^2} \right].
		\]
		we have that 
		\begin{equation}
		\label{E:Nbdd}
		N_{a,p}(u)  \le 2 \mathcal{T}_0 +  \frac{c(\sigma)}{L(\sigma)}
		\end{equation}	
		where
		\begin{align*}
		\mathcal{T}_0 = \begin{cases} 
		(\lceil \beta \rceil - 1)^{\lceil \beta \rceil -1}(ea)^{1-\lceil \beta \rceil} \Norm{u_0}_\infty  & \beta \in (0,1]
		\\ (\lceil \beta \rceil - 1)^{\lceil \beta \rceil -1}(ea)^{1-\lceil \beta \rceil} \Norm{v_0}_\infty + \Norm{u_0}_\infty & \beta \in (1,2).
		\end{cases} 
		\end{align*}
		Moreover for any $T>0$,
		\begin{equation}
		\label{E:Pnormtbdd}
		\sup_{x\in\R^d} \mathbb{E}(|u(t,x)|^p) \le \exp\left( a p t  \right) \left[ 2 \mathcal{T}_0 + \frac{ c(\sigma) }{L(\sigma)} \right]^p, \; t\in[0,T].
		\end{equation}
	\end{proposition}
	\begin{proof}
		First we fix an $a>0$ and $p \in [2,\infty)$. Using Lemma \ref{L:1and2normYZ} we see that,
		\begin{align*}
		|J(t,x)|  &= \begin{cases} 
		\left| (Z(t,\cdot) * u_0)(x)  \right| & \beta \in (0,1]
		\\	\left| (Z(t,\cdot) * v_0)(x) + (Z^*(t,\cdot) * u_0)(x) \right| & \beta \in (1,2)
		\end{cases}
		\\& \le \begin{cases}
		t^{\lceil \beta \rceil -1}\Norm{u_0}_\infty & \beta \in (0,1]
		\\ t^{\lceil \beta \rceil -1}\Norm{v_0}_\infty + \Norm{u_0}_\infty & \beta \in (1,2)
		\end{cases} 
		\end{align*}
		Now using \eqref{E:SUPidentities}, we see that 
		\begin{align}
		N_{a,p}(J) &\le \begin{cases} 
		\sup_{t\ge 0}\sup_{x\in\R}  t^{\lceil \beta \rceil -1}\Norm{u_0}_\infty e^{-at} & \beta \in (0,1]
		\\ \sup_{t\ge 0}\sup_{x\in\R} \left(  t^{\lceil \beta \rceil -1}\Norm{v_0}_\infty + \Norm{u_0}_\infty \right) e^{-at} & \beta \in (1,2)
		\end{cases} \nonumber
		\\ &= \begin{cases} 
		(\lceil \beta \rceil - 1)^{\lceil \beta \rceil -1}(ea)^{1-\lceil \beta \rceil} \Norm{u_0}_\infty  & \beta \in (0,1]
		\\ (\lceil \beta \rceil - 1)^{\lceil \beta \rceil -1}(ea)^{1-\lceil \beta \rceil} \Norm{v_0}_\infty + \Norm{u_0}_\infty & \beta \in (1,2). 
		\end{cases} \label{E:Jbdd}
		\end{align}
		
		Now we find an upper bound for $N_{a,p}(I)$. By simultaneously applying the Burkholder-Davies-Gundy inequality and triangle inequality, we get that  
		\begin{align*}
		\Norm{I(t,x)}_p^2 &\le 4p \Norm{\int_0^t \ud s \int_{\R^d}\ud y \; Y^2(t-s,x-y)\sigma^2(u(s,y))  }_{p/2}
		\\&\le 4p \int_0^t \ud s \int_{\R^d} \ud y \; Y^2(t-s,x-y)\Norm{\sigma^2(u(s,y))}_{p/2}.
		\end{align*}
		Now by applying Lemma \ref{L:1and2normYZ} and the Lipschitz property of $\sigma$ we get that 
		\begin{align*}
		\Norm{I(t,x)}_p^2 &\le 8pc^2(\sigma)C_{\#}\int_0^t (t-s)^{\rho -1} \ud s + 8pL^2(\sigma)N^2_{a,p}(u)\int_0^t \ud s \int_{\R^d} \ud y \; e^{2as} Y^2(t-s,x-y).
		\end{align*} 
		Now under Dalang's condition \eqref{E:DalangEquiv2}, the first integral above can be calculated as 
		\[
		\int_0^t (t-s)^{\rho -1} \ud s = \frac{ t^{\rho} }{\rho}.
		\]
		By applying Lemma \ref{L:1and2normYZ}, we can calculate the second integral as 
		\[
		\int_0^t \ud s \int_{\R} \ud y \; e^{2as} Y^2(t-s,x-y) = C_{\#}\int_0^t \ud s \; e^{2as}(t-s)^{\rho - 1}.
		\]
		Putting this together gives us that 
		\begin{align*}
		\Norm{I(t,x)}_p^2 &\le 8pc^2(\sigma)C_{\#} \frac{ t^{\rho} }{\rho} + 8pL^2(\sigma)N^2_{a,p}(u) C_{\#}\int_0^t \ud s \; e^{2as}(t-s)^{\rho - 1}.
		\end{align*}
		Now by multiplying both sides by $e^{-2at}$, we see that
		\begin{align*}
		\Norm{I(t,x)}_p^2e^{-2at} &\le \frac{8pc^2(\sigma)C_{\#} }{\rho} \; t^{\rho} e^{-2at} 
		\\&\qquad\qquad + 8pL^2(\sigma)N^2_{a,p}(u) C_{\#} \int_0^t \ud s \; e^{-2a(t-s)}(t-s)^{\rho - 1}.
		\end{align*}
		We now do a change of variables and again recall the definition of the Gamma function to see that 
		\begin{align*}
		\Norm{I(t,x)}_p^2e^{-2at}  &\le \frac{8pc^2(\sigma)C_{\#} }{\rho} \; t^{\rho} e^{-2at} + 8pL^2(\sigma)N^2_{a,p}(u) C_{\#} \int_0^\infty \ud u \; e^{-2au}(u)^{\rho -1} \\
		& = \frac{8pc^2(\sigma)C_{\#} }{\rho} \; t^{\rho} e^{-2at} + 8pL^2(\sigma)N^2_{a,p}(u) \frac{C_{\#}\Gamma\left(\rho \right)}{ (2a)^{\rho}}
		\end{align*}
		where the last line is true if Dalang's condition \cite{chen.hu.ea:19:nonlinear} holds. If we now apply \eqref{E:SUPidentities} and also recall that $\sqrt{a+b} \le \sqrt{a} + \sqrt{b}$, then we see that 
		\begin{align}
		N_{a,p}(I) &\le  c(\sigma)\sqrt{\frac{8pC_{\#}  }{\rho}} \left(\frac{\rho}{2ea}\right)^{\rho/2} +  L(\sigma)N_{a,p}(u) \sqrt{\frac{8pC_{\#} \Gamma\left( \rho \right)}{(2a)^\rho} }. \label{E:SoleMil3.9}
		\end{align}
		To simplify things, we now let 
		\begin{align*}
		K&:= \max\Bigg\{\sqrt{\frac{8C_{\#}  }{\rho}} \left(\frac{\rho}{2e}\right)^{\rho/2}, \sqrt{\frac{8C_{\#} \Gamma\left( \rho \right) }{2^{\rho}}}  \Bigg\}
		\end{align*}
		then,
		\begin{align}
		\label{E:Ibdd}
		N_{a,p}(I) &\le\frac{K \sqrt{p}}{a^{\rho/2}}\left( c(\sigma) + L(\sigma)N_{a,p}(u) \right).
		\end{align}
		Now by combining \eqref{E:Jbdd} and \eqref{E:Ibdd}, we see that 
		\begin{align}
		N_{a,p}(u) & \le \begin{cases} 
		(\lceil \beta \rceil - 1)^{\lceil \beta \rceil -1}(ea)^{1-\lceil \beta \rceil} \Norm{u_0}_\infty  & \beta \in (0,1]
		\\ (\lceil \beta \rceil - 1)^{\lceil \beta \rceil -1}(ea)^{1-\lceil \beta \rceil} \Norm{v_0}_\infty + \Norm{u_0}_\infty & \beta \in (1,2). 
		\end{cases} 
		\\& \quad \quad  + K\frac{\sqrt{p}c(\sigma)}{a^{\rho/2} }  + K  \frac{\sqrt{p}L(\sigma)}{a^{\rho/2}} N_{a,p}(u). \label{E:Nbddgen}
		\end{align}
		It is important to note that if we now restrict ourselves to the case of $\alpha = \beta = 2$, $d=1$ and $\gamma =0$, then \eqref{E:SoleMil3.9} and \eqref{E:Nbddgen} recover \cite[(3.9) and (3.10)]{millet.sanz-sole:21:global} respectively. Now choose $a>0$ such that 
		\[
		a > \left( 8 L^2(\sigma)K^2\right)^{1/\rho}.
		\]
		This implies that the following interval is nonempty:
		\[
		\left[ 2, \frac{a^{\rho}}{4L^2(\sigma)K^2} \right].
		\]
		In addition, for any $p$ in this interval, we have that 
		\[
		\frac{\sqrt{p}L(\sigma)K}{a^{\rho/2}} \le \frac{1}{2},
		\]
		and we also have 
		\[
		\frac{\sqrt{p}}{a^{\rho/2}}  \le \frac{1}{2KL(\sigma)}.
		\]
		This along with \eqref{E:Nbddgen} implies that 
		\begin{align*}
		N_{a,p}(u) & \le 2 \begin{cases} 
		(\lceil \beta \rceil - 1)^{\lceil \beta \rceil -1}(ea)^{1-\lceil \beta \rceil} \Norm{u_0}_\infty  & \beta \in (0,1]
		\\ (\lceil \beta \rceil - 1)^{\lceil \beta \rceil -1}(ea)^{1-\lceil \beta \rceil} \Norm{v_0}_\infty + \Norm{u_0}_\infty & \beta \in (1,2)
		\end{cases} 
		\\& \quad \quad + \frac{c(\sigma)}{L(\sigma)}
		\end{align*}
		and this is precisely \eqref{E:Nbdd}. Note that \eqref{E:Pnormtbdd} follows immediately from \eqref{E:Nbdd}.  
	\end{proof}
	
	We will need to recall the following result proved in \cite{chen.hu.ea:19:nonlinear} in order to prove Proposition \ref{P:NormInc} below. 
	
	\begin{proposition}{\normalfont\cite[Proposition 2.1 and 2.2]{chen.hu:21:holder}}
		\label{P:YInc}
		Suppose that $\alpha \in (0,2]$, $\beta \in (0,2)$, $\gamma>0$, and \eqref{E:DalangEquiv} holds. Consider $\rho = \rho(d)$ which is defined above in \eqref{D:rho}. Then $Y(t,x) = Y_{\alpha, \beta, \gamma, d}(t,x)$ satisfies the following:
		\begin{enumerate}
			\item \label{P:YIncCase1}  Suppose that $0< \theta < (\Theta -d) \wedge 2$ where $\Theta$ is given in \eqref{E:DalangEquiv}, $0< t,r \le T$ for some $T>0$, $\beta \le 1$ and $\gamma \le \lceil \beta \rceil - \beta$. Then there exists some $C := C(\alpha, \beta, \gamma, \nu, d, \theta, T)$ such that
			\[
			\int_0^T \ud s \int_{\R^d} \ud y |Y(t-s,x-y) - Y(r-s,z-y)|^2 \le C \left(|t-r|^\rho + |x-z|^{\theta}\right).
			\]
			\item \label{P:YIncCase2} If we only assume that $\alpha \in (0,2]$, $\beta \in (0,2)$, $\gamma>0$, $0< t,r \le T$ for some $T>0$ and \eqref{E:DalangEquiv} holds then for $0< \theta < (\Theta -d) \wedge 2$ there exists some $C := C(\alpha, \beta, \gamma, \nu, d, \theta, T)$ such that
			\[
			\int_0^T \ud s \int_{\R^d} \ud y |Y(t-s,x-y) - Y(r-s,z-y)|^2 \le C \left(|t-r|^q + |x-z|^{\theta}\right),
			\]
			where 
			\[
			0 < q < \rho.
			\]
		\end{enumerate}
	\end{proposition}
	
	\begin{remark}
		In Proposition \ref{P:YInc} above, we may often ignore the dependence of $C$ on $\alpha, \beta, \gamma, \nu, d$ and just write $C(\theta, T)$. 
	\end{remark}
	
	\begin{proposition}
		\label{P:NormInc}
		Let $a>0$ and suppose that $0< \theta < (\Theta -d) \wedge 2$ and $\rho = \rho(d) > 0$. Then for any $x,z \in \R$, $p\ge 2$ and for any $T>0$ with $t\in [0,T]$, the stochastic integral, $I(t,x)$, satisfies the following:
		\begin{equation}
		\label{E:NormInc_N(u)}
		\frac{\Norm{I(t,x)-I(r,z)}_p}{\left(|t-r|^{q} + |x-z|^{\theta}\right)^{1/2}} \le 	C(p,\theta, T)\left[ \mathcal{M}_1 + \mathcal{M}_2 e^{aT} N_{a,p}(u)\right] ,
		\end{equation}
		where 
		\[
		\mathcal{M}_1 = \sqrt{p}c(\sigma) \quad \text{and} \quad\mathcal{M}_2 =\sqrt{p}L(\sigma),
		\]
		and
		\begin{enumerate}
			\item $0< q \le \rho$ under Case \ref{P:YIncCase1} of Proposition \ref{P:YInc},
			\item $0 < q < \rho$ under Case \ref{P:YIncCase2} of Proposition \ref{P:YInc}.
		\end{enumerate}
		Moreover, for any 
		\[
		a > \left( 8 L^2(\sigma)K^2\right)^{1/
			\rho}.
		\]
		and 
		\[
		p \in \left[ 2, \frac{a^{\rho}}{4L^2(\sigma)K^2} \right]
		\]
		we have that 
		\begin{equation}
		\label{E:NormInc_cL}
		\frac{\Norm{I(t,x)-I(r,z)}_p}{\left(|t-r|^{q} + |x-z|^{\theta}\right)^{1/2}} \le C(p,\theta, T)\left[ \mathcal{M}_1 + \mathcal{M}_2 e^{aT} \left( 2 +  \frac{c(\sigma)}{L(\sigma)} \right)\right].
		\end{equation}
	\end{proposition}
	
	\begin{proof}
		Note that because of the initial conditions, we see that $\Norm{u(t,x) - u(r,z)}_p = \Norm{I(t,x) - I(r,z)}_p$. We now apply Burkholder-Davies-Gundy inequality along with the triangle inequality, as was done in the proof of Proposition \ref{P:N(u)Bdd},  to see that
		\begin{align*}
		\Norm{I(t,x) - I(r,z)}_p^2 &\le 4p \Norm{ \int_0^T \ud s \int_{\R^d} \ud y \left| Y(t-s,x-y) - Y(r-s,z-y) \right|^2\sigma^2(u(s,y)) }_{p/2} \\
		&\le 4p \int_0^T \ud s \int_{\R^d} \ud y \left| Y(t-s,x-y) - Y(r-s,z-y) \right|^2  \Norm{\sigma^2(u(s,y)) }_{p/2}.
		\end{align*} 
		Now by applying the Lipchitz condition on $\sigma$ gives us that 
		\begin{align*}
		\Norm{I(t,x) - I(r,z)}_p^2 &\le 8p \int_0^T \ud s \int_{\R^d} \ud y \left| Y(t-s,x-y) - Y(r-s,z-y) \right|^2 \left( c^2(\sigma) + L^2(\sigma)\Norm{u(s,y)}_{p}^2\right) \\
		&\le 8p\Bigg[ c^2(\sigma)   \int_0^T \ud s \int_{\R^d} \ud y \left| Y(t-s,x-y) - Y(r-s,z-y) \right|^2 \\
		& \qquad+ L^2(\sigma) \int_0^T \ud s \int_{\R^d} \ud y \left| Y(t-s,x-y) - Y(r-s,z-y) \right|^2 \exp\left(2as\right)N_{a,p}^2(u)  \Bigg] \\
		&\le 8p\Bigg[ c^2(\sigma)   \int_0^T \ud s \int_{\R^d} \ud y \left| Y(t-s,x-y) - Y(r-s,z-y) \right|^2 \\
		& \qquad+ L^2(\sigma) \exp\left(2aT \right)N_{a,p}^2(u)  \int_0^T \ud s \int_{\R^d} \ud y \left| Y(t-s,x-y) - Y(r-s,z-y) \right|^2  \Bigg]. 
		\end{align*}
		We now apply Proposition \ref{P:YInc} to see that 
		\begin{align*}
		\Norm{I(t,x) - I(r,z)}_p^2 &\le 8p \bigg[ c^2(\sigma)C(\theta, T)\left(|t-r|^{q} + |x-z|^{\theta}\right)\\ 
		& \qquad +  L^2(\sigma) \exp\left(2aT \right)N_{a,p}^2(u) C(\theta, T)\left(|t-r|^{q} + |x-z|^{\theta}\right) \bigg],
		\end{align*}
		where $C(\theta, T)$ is as in Proposition \ref{P:YInc}. Now by taking the square root of both sides and applying the identity $\sqrt{x+y} \le \sqrt{x} + \sqrt{y}$ yields
		\begin{align*}
		\Norm{I(t,x) - I(r,z)}_p &\le 2\sqrt{2C(\theta, T)p} \Big[ c(\sigma) + L(\sigma) \exp\left(aT \right)N_{a,p}(u) \Big]\left(|t-r|^{q} + |x-z|^{\theta}\right)^{1/2}. 
		\end{align*}
		This proves \eqref{E:NormInc_N(u)}. It is clear that \eqref{E:NormInc_cL} follows directly from \eqref{E:NormInc_N(u)} and Proposition \ref{P:N(u)Bdd}. 
	\end{proof}
	
	\begin{proposition} \label{P:Holder}
		The stochastic integral, $I$, has a version, still denoted by $I$ that is $\eta_1$-H\"older continuous in time and $\eta_2$-H\"older continuous in space with $0 < \eta_1 < q/2$ and $0 < \eta_2 < \theta/2$ where $0 < \theta < (\Theta -d) \wedge 2$ and 
		\begin{enumerate}
			\item $0< q \le \rho$ under Case \ref{P:YIncCase1} of Proposition \ref{P:YInc},
			\item $0 < q < \rho$ under Case \ref{P:YIncCase2} of Proposition \ref{P:YInc}.
		\end{enumerate}
		Moreover, if we consider
		\[
		a > \left( 8 L^2(\sigma)K^2\right)^{1/\rho} \quad \text{and} \quad 	p \in \left[ 2, \frac{a^{\rho}}{4L^2(\sigma)K^2} \right],
		\]
		where $K$ is the constant from Proposition \ref{P:N(u)Bdd}, then  
		\[
		\mathbb{E}\left[ \sup_{t\in[0,T],\; |x| \le R} |u(t,x)|^p \right] \le 2^{p-1}\widehat{C}_T^p  + C(p, \theta, T,R) \left[ \mathcal{M}_1^p + \mathcal{M}_2^p e^{apT} \left(2 +  \frac{c(\sigma)}{L(\sigma)}\right)^p\right],
		\]
		where $\widehat{C}_T$ is defined in \eqref{E:HomogBdd}.
	\end{proposition}
	
	\begin{proof}
		Proposition \ref{P:NormInc} implies that 
		\begin{equation}
		\Norm{I(t,x) - I(r,z)}_p \le C(p,\theta, T)\left[ \mathcal{M}_1 + \mathcal{M}_2e^{aT}N_{a,p}(u) \right]\left(|t-r|^{q} + |x-z|^{\theta}\right)^{1/2}.
		\end{equation}
		By Kolmogorov's continuity theorem, there exists a version of $I$, which we denote by $I$, that is $\eta_1$-H\"older continuous in time and $\eta_2$-H\"older continuous in space with with $0<\eta_1 < q/2$ and $0<\eta_2 < \theta/2$.
		
		Next, note that by triangle inequality,
		\begin{align*}
		|u(t,x)|^p &\le 2^{p-1}\left( \Norm{J_0(\cdot,\circ)}_{\infty}^p + |u(t,x)-J_0(t,x)|^p \right) \\
		&\le 2^{p-1}\left( \widehat{C}_T^p + \sup_{t\in[0,T],\; |x| \le R} \: |I(t,x)|^p \right) \\
		&= 2^{p-1}\left( \widehat{C}_T^p + \sup_{t\in[0,T],\; |x| \le R} \: |I(t,x)|^p\right),
		\end{align*}
		where we recall that $\widehat{C}_T$ defined in \eqref{E:HomogBdd}. Due to the continuity of $I$ on the compact set $[0,T] \times \bar{B}_R(0)$, there exists a point $(t_0,x_0) \in [0,T] \times \bar{B}_R(0)$ such that 
		\[
		\sup_{t\in[0,T],\; |x| \le R} \: |I(t,x)|^p = |I(t_0,x_0)|^p.
		\]
		Thus, for any $(t,x) \in [0,T] \times \bar{B}_R(0)$,
		\[
		|u(t,x)|^p \le 2^{p-1}\left( \widehat{C}_T^p  + |I(t_0,x_0)|^p \right)
		\] 
		which in return implies that 
		\[
		\sup_{t\in[0,T],\; |x| \le R} |u(t,x)|^p \le 2^{p-1}\left(\widehat{C}_T^p + |I(t_0,x_0)|^p \right).
		\]
		Now by taking the expectation of both sides and applying Proposition \ref{P:NormInc} gives us 
		\begin{align*}
		\mathbb{E}\left[ \sup_{t\in[0,T],\; |x| \le R} |u(t,x)|^p \right] &\le  2^{p-1}\widehat{C}_T^p  + 2^{p-1}C(p, \theta, T)^p\left((2T)^{q/2} + (2R)^{\theta /2}\right)^p\left[ \mathcal{M}_1 + \mathcal{M}_2 e^{aT} N_{a,p}(u)\right]^p \\
		&= 2^{p-1}\widehat{C}_T^p  + C(p, \theta, T,R) \left[ \mathcal{M}_1^p + \mathcal{M}_2^p e^{apT} N_{a,p}(u)^p\right].
		\end{align*}
		Now if we consider
		\[
		a > \left( 8 L^2(\sigma)K^2\right)^{1/\rho} \quad \text{and} \quad 	p \in \left[ 2, \frac{a^{\rho}}{4L^2(\sigma)K^2} \right],
		\]
		then by Proposition \ref{P:N(u)Bdd}, we have that 
		\[
		\mathbb{E}\left[ \sup_{t\in[0,T],\; |x| \le R} |u(t,x)|^p \right] \le 2^{p-1}\widehat{C}_T^p  + C(p, \theta, T,R) \left[ \mathcal{M}_1^p + \mathcal{M}_2^p e^{apT} \left(2 +  \frac{c(\sigma)}{L(\sigma)}\right)^p\right].
		\]
	\end{proof}
	
	\begin{remark}
		The H\"older continuity in time proven above in Proposition \ref{P:Holder} recovers the result proven by Foondun and Nane in \cite[Theorem 1.9]{foondun.nane:17:asymptotic} if we restrict ourself to the case where $\beta \in (0,1)$ and $\gamma = 1-\beta$. 
	\end{remark}

	
	\section{Proof of Theorem \ref{T:MainSupLin}} \label{S:TMainSupLin}
	
	\begin{proof}
		We first solve a {\it truncated} version of \eqref{E:SupLinSPDE}. By this, we mean $\eqref{E:SupLinSPDE}$ with $\sigma$ replaced by $\sigma_N$ where in the case of $d=1$, 
		\[
		\sigma_N(x) = \sigma(x)\textbf{1}_{\{|x| \le N\}} + \sigma(N)\textbf{1}_{\{x  >N\}} + \sigma(-N)\textbf{1}_{\{x \le -N\}},
		\]
		and in higher dimensions $\sigma_N$ is defined similarly. Note that since $\sigma$ satisfies \eqref{C:SigSupLin}, then $\sigma_N$ is Lipschitz continuous with Lipschitz constant $L(\sigma_N) := \sigma_2 \ln(2N)^{\delta}$. In other words, 
		\[
		|\sigma_N(x)| \le \sigma_1 + \sigma_2 \ln(2N)^{\delta}|x|.
		\]
		Hence we can apply Theorem \ref{T:SolveabilityGenLip} and see that there exists an unique solution of the truncated version of \eqref{E:SupLinSPDE}, which we denote by $u_N := \left\{ u_N(t,x) : (t,x) \in [0,T] \times \R \right\}$. In addition, Proposition \ref{P:Holder} implies that the solution has a version, still denoted by $u_N$, which is $\eta_1$-H\"older continuous in time and $\eta_2$-H\"older continuous in space  with $0<\eta_1 < q/2$ and $0<\eta_2 < \theta/2$.
		
		We will now apply Proposition \ref{P:Holder} to find an upper bound on the $p$-norm of $u_N$. For this, consider 
		\[
		a > \left( 8 L^2(\sigma_N) K^2 \right)^{1/\rho} \quad \text{and} \quad p \in \left[ 2, \frac{a^{ \rho}}{4L^2(\sigma_N)K^2} \right].
		\] 
		Then Proposition \ref{P:Holder} implies that 
		\begin{equation}
		\label{E:u_NMomBdd}
		\mathbb{E}\left(\sup_{t\in[0,T],\; x\in[-R,R]} |u_N(t,x)|^p\right) \le 2^{p-1} + C(p, \theta, T,R)\left[\mathcal{M}_1^p + \mathcal{M}_2^p(N) \exp(apT)\left(2+ \frac{\sigma_1}{L(\sigma_N)}\right)\right]
		\end{equation}
		where we recall that 
		\begin{equation}
		\label{D:M1&2forSig_N}
		\mathcal{M}_1 = \sqrt{p}\sigma_1 \quad \text{and} \quad \mathcal{M}_2(N) = \sqrt{pL(\sigma_N)} = \sqrt{p\sigma_2\ln(2N)^\delta}.
		\end{equation}
		We now use the above to prove the existence and uniqueness of \eqref{E:SupLinSPDE} with superlinear $\sigma$ satisfying \eqref{C:SigSupLin}. For any $T>0$ and $N \ge 2$ with $N\in \mathbb{N}$, we define the following stopping time:
		\begin{equation}
		\tau_N := \inf \left\{ t>0 : \sup_{|x| \le R} |u_N(t,x)| \ge N \right\} \wedge T.
		\end{equation}
		The uniqueness of the solution and the local property of the stochastic integral imply that on $\{t < \tau_N\}$, $u_N(t,x) = u_{N+k}(t,x)$ for any $k \in \mathbb{N}$. Hence, $(\tau_N)_{N \ge 2}$ is increasing and bounded above by $T$. 
		
		We now momentarily assume that $\sup_N \tau_N = T$ which would imply that $\{t < \tau_N \} \uparrow \Omega$. On each $\{ t < \tau_N \}$, define $(u(t,x) : (t,x) \in [0,T) \times \R)$ by $u(t,x) = u_N(t,x)$ and hence $u(t,x) = u_{N+k}(t,x)$ for any $k\in \mathbb{N}$. This implies that on $\{ t < \tau_N \}$ 
		\[
		u(t,x) = 1 + \int_0^t \ud s \int_{\R} \ud y \; Y(t-s,x-y) \sigma_N(u(s,y)) W(\ud s, \ud y).
		\]
		However, on $\{ t < \tau_N \}$, we have that $\sigma_N(u_N(t,x)) = \sigma(u(t,x))$ and so $u$ satisfies
		\[
		u(t,x) = 1 + \int_0^t \ud s \int_{\R} \ud y \; Y(t-s,x-y) \sigma(u(s,y)) W(\ud s, \ud y), \quad \{ t < \tau_N \}.
		\]
		Lastly, since $\{ t < \tau_N \} \uparrow \Omega$, we conclude that 
		\begin{equation}
		\label{E:SupLinMildSol}
		u(t,x) = 1 + \int_0^t \ud s \int_{\R} \ud y \; Y(t-s,x-y) \sigma(u(s,y)) W(\ud s, \ud y), \quad (t,x) \in [0,T) \times \R.
		\end{equation}
		It is mentioned that the stochastic integrals are no longer defined in $L^2(\Omega)$. The extension of the integral is found in Dalang and Sanz-Sole's textbook that is not yet published.
		
		We now focus on proving that $\sup_N \tau_N = T$, or equivalently that $ \mathbb{P}\left(\tau_N <T\right) \to 0$ as $N \to \infty$. To do this, note that by Chebychev's inequality
		\begin{align*}
		\mathbb{P}\left( \tau_N < T \right) &\le \mathbb{P}\left(\sup_{t\in[0,T],\; |x| \le R} |u_N(t,x)| \ge N \right) \\
		& \le N^{-p}\mathbb{E}\left( \sup_{t\in[0,T],\; |x| \le R} |u_N(t,x)|^p \right).
		\end{align*}
		An application of \eqref{E:u_NMomBdd} gives 
		\begin{align*}
		N^{-p}\mathbb{E}\left( \sup_{t\in[0,T],\; |x| \le R} |u_N(t,x)|^p \right) &\le N^{-p}\left(  2^{p-1} + C(p, \theta, T,R) \left[ \mathcal{M}_1^p + \mathcal{M}_2(N)^p \exp\{apT\} \left(2 +  \frac{\sigma_1}{L(\sigma_N)}\right)^p\right]\right) \\
		&= \frac{2^{p-1} + C(p, \theta, T,R)\mathcal{M}_1^p }{N^p} + \frac{C(p,T,R) \mathcal{M}_2(N)^p \exp\{apT\} \left(2 +  \frac{\sigma_1}{L(\sigma_N)}\right)^p }{ N^p }.
		\end{align*}
		The first term clearly converges to $0$ as $N\to \infty$. As for the second term, we expand using \eqref{D:M1&2forSig_N} to see that 
		\begin{align*}
		\frac{\mathcal{M}_2(N)^p \exp\{apT\} \left(2 +  \frac{\sigma_1}{L(\sigma_N)}\right)^p }{ N^p } &= \frac{\left(\sqrt{p\sigma_2\ln(2N)^\delta}\right)^p \exp\{apT\} \left(2 +  \frac{\sigma_1}{ \sqrt{p\sigma_2\ln(2N)^\delta} }\right)^p }{ N^p }  \\
		&\sim \frac{ \ln(2N)^{p\delta/2} }{ N^p } \to 0,
		\end{align*}
		where above we use the symbol $\sim$ to mean asymptotically equivalent as $N \to \infty$. This completes the proof of the Theorem. 
	\end{proof}

\section{Some comments on a spatially colored noise}
Consider equation \ref{E:SupLinSPDE} driven by a noise that is white in time and colored space. In other words, $W$ is a $0$-mean Gaussian process indexed on the space of Schwartz functions with the following covariance structure:
	\begin{equation} \label{E:Colored} 
	C(\phi, \psi) := \int_0^\infty \ud t \int_{\R^{d}}  \Gamma(\ud x) \; (\phi(t,\cdot)*\tilde{\psi}(t,\cdot))(x), \quad \tilde{\psi}(x) = \psi(-x).
	\end{equation}
As per usual, the tempered measure $\Gamma$ is non-negative and nonnegative definite and satisfies the following two conditions:
	\begin{assumption} \label{A:NoiseAssum1} $\;$
	\begin{enumerate}
		\item the spectral measure $\mu$ is such that $\int_{\R^d} \frac{\mu(\ud \zeta)}{1 + |\zeta|^\alpha} < \infty$,
		\item $\Gamma$ is absolutely continuous with respect to the Lebesgue measure and therefore has a density, $\Gamma(\ud x) = f(x) \ud x$. Furthermore, when this is the case, we may often write $\mu(\ud \zeta) = \varphi(\zeta) \ud \zeta$ where we have the relation $\varphi = \mathcal{F}f$,
	\end{enumerate}
	\end{assumption}

\noindent Under this assumption, we may rewrite \eqref{E:Colored} as 
\[
C(\phi,\psi) = \int_0^\infty \ud t \int_{\R^{2d}} \ud x \ud y \; \phi(t,x) f(x-y) \psi(t,y)
\]
and by applying Placherel's identity, this further implies that 
\[
\int_0^\infty \ud t \int_{\R^{2d}} \ud x \ud y \; \phi(x) f(x-y) \psi(y) = (2 \pi)^{-d} \int_0^\infty \ud t \int_{\R^{2d}} \mu(\ud \zeta) \mathcal{F}(\phi(t,\cdot))(\zeta) \overline{\mathcal{F}(\psi(t,\cdot))(\zeta)}.
\]

We now prove a lemma which is a generalization of equation \eqref{C:Dalang} which accounts for the colored noise.

	\begin{lemma}
	\label{L:YfYbdd}
	Suppose Assumption \ref{A:NoiseAssum1} holds. Then there exists a constant $C(\beta,\gamma) =: C >0$ such that for any $x,y \in \R^d$ we have that 
	\begin{equation}
	\label{E:YYfBdd}
	\int_{\R^{2d}} Y(t,x-z)Y(t,y-w)f(w-z) \; \ud z \ud w \le C \left( \frac{\nu}{2} \right)^{-\lambda / \alpha} t^{2(\beta + \gamma -1)-\beta \lambda / \alpha} \int_{\R^{d}}  \frac{ \mu(\ud \zeta) }{\left(1+ |\xi|^\alpha\right)^2 } < \infty
	\end{equation}
\end{lemma}
\begin{proof}
	Through a change of variables, we may write the left side of \eqref{E:YYfBdd} as the following:
	\[
	\int_{\R^{2d}} Y(t,x-z)Y(t,y-w)f(w-z) \; \ud z \ud w = \int_{\R^{d}} Y\left(t,u-(x-y)\right)(Y(t,\cdot)*f)(u) \; \ud u,
	\]
	where $*$ denotes the spatial convolution. Now an application of Pareseval's identity gives us that 
	\[
	\int_{\R^{d}} Y\left(t,u-(x-y)\right)(Y(t,\cdot)*f)(u) \; \ud u = \frac{1}{(2\pi)^d} 	\int_{\R^{d}} \mathcal{F}Y\left(t,\cdot-(x-y)\right)(\zeta) \mathcal{F}(Y(t,\cdot)*f)(\zeta) \ud \zeta.
	\]
	Applying the translation property of the Fourier transform and \cite[Theorem 4.1]{chen.hu.ea:19:nonlinear} gives us that 
	\begin{align*}
	\frac{1}{(2\pi)^d} 	\int_{\R^{d}} \mathcal{F}Y\left(t,\cdot-(x-y)\right)(\zeta) & \mathcal{F}(Y(t,\cdot)*f)(\zeta) \ud \xi \\
	& = \frac{ t^{2(\beta + \gamma -1)}}{ (2\pi)^d } \int_{\R^{d}} \exp(-i(x-y)\zeta) E_{\beta,\beta+\gamma}\left(-\frac{\nu t^\beta |\zeta|^\alpha}{2}\right)^2 \mu(\ud \zeta).
	\end{align*}
	Hence by taking the absoulte value on both sides and recalling the positivity of $f$, $Y$ and the Mittag-Leffler function we get that 
	\begin{align*}
	\int_{\R^{2d}} Y(t,x-z)Y(t,y-w)f(w-z) \; \ud z \ud w \le \frac{ t^{2(\beta + \gamma -1)} }{ (2\pi)^d } \int_{\R^{d}}  E_{\beta,\beta+\gamma}\left(-\frac{\nu t^\beta |\zeta|^\alpha}{2}\right)^2 \mu(\ud \zeta).
	\end{align*}
	Now, by applying \cite[Theorem 1.6]{podlubny:99:fractional}, we see that there exists a $C>0$ depending only on $\beta$ and $\gamma$ such that 
	\[
	\frac{ t^{2(\beta + \gamma -1)}}{ (2\pi)^d } \int_{\R^{d}}  E_{\beta,\beta+\gamma}\left(-\frac{\nu t^\beta |\xi|^\alpha}{2}\right)^2 \mu(\ud \zeta).\le 
	\frac{C  t^{2(\beta + \gamma -1)}}{ (2\pi)^d } \int_{\R^{d}}  \frac{ \mu(\ud \zeta) }{\left(1+ \frac{\nu}{2} t^\beta |\zeta|^\alpha \right)^2 } .
	\]
	Another change of variables implies that 
	\[
	\frac{C  t^{2(\beta + \gamma -1)}}{ (2\pi)^d } \int_{\R^{d}}  \frac{ \mu(\ud \zeta) }{\left(1+ \frac{\nu}{2} t^\beta |\zeta|^\alpha \right)^2 }  \le 
	\frac{C }{ (2\pi)^d } \left(\frac{\nu}{2}\right)^{-\lambda/\alpha}  t^{2(\beta + \gamma -1)-\beta \lambda/\alpha} \int_{\R^{d}}  \frac{ \mu(\ud \zeta) }{\left(1+ |\xi|^\alpha\right)^2 } .
	\]
	The finiteness of the above integral follows from Assumption \ref{A:NoiseAssum1}.
\end{proof}	

	\addcontentsline{toc}{section}{References}
	\printbibliography[title={References}]
%\begin{thebibliography}{999}
%	%sample
%	%\bibitem{Chen17Tran}% {{{ 3.
%	%Chen, L. (2017).
%	%\newblock Nonlinear stochastic time-fractional diffusion equations on $\R$: moments, H\"older regularity and intermittency.
%	%\newblock {\em Trans. Amer. Math. Soc.} \textbf{369} (12) 8497--8535.
%	\bibitem{CH2017}
%	Chen L. and Huang J. (2017).
%	\newblock Comparison Principle for SHE on $\R^d$.
%	\newblock {\em The Annals of Probability} \textbf{Vol. 47, No. 2} 989 -- 1035.
%	
%	\bibitem{ChenHu21}
%	Chen, L and Hu, G. (2021).
%	\newblock H\"older regularity of the nonlinear stochastic time-fractional slow and fast diffusion equations on $\R^d$.
%	\newblock{\em arXiv:2105.00891.}
%	
%	\bibitem{CHHH16}
%	Chen, L., Hu, G., Hu, Y. and Huang, J. (2016)
%	\newblock{Space-time fractional diffusions in Gaussian noisy environment}
%	\newblock{\it Stochastics}, 89:1, 171--206.
%	
%	\bibitem{CHN19}% {{{ 5.
%	Chen, L., Hu, Y., and Nualart, D. (2019).
%	\newblock Nonlinear stochastic time-fractional slow and fast diffusion equations on $\R^d$.
%	\newblock {\em Stochastic Process. Appl.} \textbf{129}, 5073--5112
%	% }}}
%	
%	\bibitem{FN2016}
%	Foondun, M and Nane E. (2016)
%	\newblock{Asymptotic properties of some space-time fractional stochastic equations.}
%	\newblock{\em Mathematische Zeitschrift} 2015, \textbf{287}, 493 -- 519
%	
%	\bibitem{FT2014}
%	Foondun, M. and Tian, K. (2014)
%	\newblock On some properties of a class of fractional stochastic heat equations.
%	\newblock {\em arXiv:1404.6791v1}
%	
%	\bibitem{HHN13}
%	Hu, Y., Huang, J. and Nualart, D. (2013)
%	\newblock On H\"older continuity of the solution of stochastic wave
%	equations
%	\newblock {\it arXiv:1308.6607v1}
%	
%	\bibitem{Kilbas}
%	\newblock H-Transforms: Theory and Applications 
%	\newblock{\it An International Series of Monographs in Mathematics.}
%	
%	\bibitem{Kuo}
%	Kuo, H. (2000)
%	\newblock{\em Introduction to Stochastic Integration}
%	\newblock Springer, Universitext
%	
%	\bibitem{Podlubny99}
%	Podlubny, I. (1999).
%	\newblock{\em Fractional Differential Equations.}
%	\newblock Mathematics in Science and Engineering, Academic Press, Vol198 (1999).
%	
%	\bibitem{SoleMillet2021}
%	Sanz-Sole, M. and Millet, A. (2021)
%	\newblock{Global solutions to stochastic wave equations with superlinear coefficients.}
%	\newblock {\em Stochastic Processes and their Applications}, Elsevier, 2021, 139, pp.175-211.
%	
%	\bibitem{Walsh89}
%	Walsh, J. (1989)
%	\newblock{\em An Introduction to Stochastic Partial Differential Equations.}
%	\newblock \'Ecole d'\'et\'e de Probabilit\'es de Saint-Flour, XIV|1994. Lecture Notes in Math., vol. 1180, pp. 265-439. Springer, Berlin (1986).
%\end{thebibliography}

\end{document}